\section{Introduction}
\label{sec:intro}

How long should a task instruction be?
Prompts used in practice have grown steadily longer: system prompts for assistants and agents now often run to hundreds or thousands of words, full of rules, edge cases, rationale, and reminders.
The usual intuition is that more detail can only help, since the model gets more information about what we want.
A competing intuition, often stated as ``attention dilution,'' holds that long prompts bury the key requirement and make models worse.
Both views shape how people write prompts, yet neither has been tested in a controlled way.

The question matters most for the settings where prompts are cheapest to change and models are weakest.
Small open-weight models run on-device or at low cost, and practitioners often copy long instructions written for frontier models into them.
If long instructions hurt, and if the harm depends on task type or model size, then prompt length is a design choice with real consequences rather than a matter of taste.

Prior work gives indirect evidence.
Models degrade when the \emph{input context} grows, even with irrelevant text~\citep{shi2023distracted,liu2024lost,levy2024same}, and instruction following degrades as the \emph{number of constraints} grows~\citep{zhou2023ifeval,jiang2024followbench}.
However, these studies vary the input or the constraint set, not the length and detail of the task instruction itself.
They also rarely separate two explanations for a length effect: \emph{more tokens} (dilution) \versus \emph{more content} (over-specification, interference from extra rules).

{\bf What do we do?}
We build a nested ladder of five instruction tiers (\minimal, \concise, \standard, \detailed, \exhaustive) for each of four task types: classification, extraction, generation, and reasoning.
Each tier adds one block of content to the previous tier, so a longer tier never removes information, and every tier states the same core output-format constraint.
We add a \padded control that has the \standard content plus task-irrelevant filler, word-matched to \exhaustive.
We evaluate five instruction-tuned \qwen models (0.5B, 1.5B, 3B, 7B, 14B) on 252 fixed items per task, and we score task correctness and format compliance separately.
We pre-registered six hypotheses, including that classification and extraction peak at \standard and drop 3--7 points at \exhaustive, and that 7B models suffer more than larger ones.

{\bf What do we find?}
Most hypotheses fail.
For 3B--14B models, the \standard$\to$\exhaustive change is within $\pm$1.6 points on every task.
After Holm correction over the 20 primary tests, the only significant change is a $+1.0$ ROUGE-L \emph{gain} for 7B on summarization.
The 7B and 14B models respond to length in the same way (difference-in-differences between $-1.4$ and $+1.2$ points, all n.s.).
Harm is confined to the smallest models: the 0.5B model's classification accuracy drops from 62.3\% at \minimal to 43.7\% at \standard and 39.3\% at \exhaustive, as its predictions collapse onto ``Business.''
The size of the \standard$\to$\exhaustive drop shrinks with model size (Spearman $\rho=-0.9$ for both classification and extraction).
The padding control shows that token count is not the cause: filler never significantly hurts any model, and it lifts 0.5B accuracy from 43.7\% to 59.5\%.
The largest single gain from instruction detail comes earlier in the ladder: going from \minimal to \concise raises 7B and 14B extraction F1 by about 10 points.

Our contributions are:
\begin{itemize}[leftmargin=*,itemsep=0pt,topsep=0pt]
    \item We design a controlled, nested five-tier instruction ladder for four task types, with a fixed format constraint in every tier and a length-matched padding control that separates token count from content.
    \item We conduct a paired evaluation of five \qwen models (0.5B--14B) over 6 conditions and 4 tasks (30{,}240 generations), with bootstrap confidence intervals, permutation tests, and Holm correction.
    \item We show that instruction length is roughly neutral for models of 3B parameters and above, and that harm from long instructions is a small-model effect caused by instruction content and its position, not by length.
    \item We give concrete prompt-writing guidance: include a short task definition, restate the output format at the end of long prompts, and keep instructions minimal (or in the system message) for sub-2B models.
\end{itemize}

\Secref{sec:related} reviews related work, \secref{sec:method} describes the design, \secref{sec:results} reports results, \secref{sec:discussion} interprets them and lists limitations, and \secref{sec:conclusion} concludes.
